\section{Discussion}
\label{sec:discussion}

\paragraph{What the prespecified tests show}
Under a margin fixed before the curve was computed, four of the 16 GRABMyo electrodes could be removed while remaining non-inferior to the full array. Five could not. Every $k$ from 12 to 15 met the margin, so the result does not rest on a single \rev{isolated} crossing\rev{, and the plan's Holm correction, applied post hoc, left $k^*$ unchanged}. The margin is an engineering tolerance, not a clinical threshold. Every reduction cost some Macro-F1, and halving the array cost about \DeltaKEight. The practical reading is modest: in a subject-independent setting, this 16-electrode \rev{monopolar} array contains only limited redundancy for 16-gesture recognition. User-dependent studies that retain far fewer channels answer a different question~\cite{phinyomark2018emg,qu2021reduce}, and their results should not be taken as evidence that a \rev{layout intended for many users} can be equally sparse.

Two of the three hypotheses about attribution were supported. Attribution-ranked subsets \rev{exceeded the mean of the random subsets} at every size tested (H2), and deletion with retraining showed the expected ordering at every size (H4). \rev{Deletion with retraining therefore supports that the ranking reflects channels on which the retrained models depend.} The stability penalty, however, \rev{produced no measurable benefit} (H3). With cross-fitting over roughly 34 participants per inner-training set, inner-fold attributions already agreed so closely that the dispersion term was a few percent of the magnitude term. \rev{It changed the top-12 set in at most one fold.} We report SR-CFA as a tested and unsupported variant. \rev{The evidence therefore favours plain cross-fitted TreeSHAP magnitude over the regularised variant;} a related robustness-weighted score has been proposed independently~\cite{xiang2026roshap}.

\paragraph{\rev{Limited additional benefit over simple alternatives}}
The exploratory comparisons temper H2. \rev{Exceeding the mean of random subsets} is a low bar. \rev{In exploratory comparisons, neither a data-free geometric rule nor permutation importance differed detectably from the attribution ranking at 4, 8 and 12 electrodes; attribution was higher at $k = 6$. Absence of a detectable difference is not evidence of equivalence.} Within each participant's own random subsets, the attribution subset was typically above the median but rarely near the top. \rev{One hypothesis consistent with this pattern is that even circumferential coverage accounts for much of what a ranking achieves. Whether a data-free, evenly spread layout is itself non-inferior to the full array was not tested and is a question for future work.}

\paragraph{Interpretation of the ranked electrodes}
The four highest-ranked electrodes (\TopFourSet) were the same in \TopFourFolds\ of 43 folds and in single-participant rankings (mean Spearman $\rho$ = \StabSubjRho). We deliberately do not map them to named muscles. The dataset documents ring positions, not the anatomical placement of each circumferential index\rev{, and no imaging-based registration of electrodes to muscles was available}. Surface electrodes also record compound activity from several muscles at each site~\cite{farina2004extraction}. The robust statement is positional. High attribution concentrated at a few circumferential positions (Fig.~\ref{fig:attribution}a). The top and bottom of the ranking were stable, and the ordering varied mainly among mid-ranked electrodes. \rev{Physiologically, it is plausible that a few circumferential sites capture activity shared by several gestures, but these data cannot attribute that activity to specific muscles.}

\paragraph{Cross-session and replication}
Cross-session results were exploratory and should be read narrowly. \rev{No additional loss on later days was detected on top of the cross-user gap.} This \rev{does not show} that electrode shift is harmless~\cite{hargrove2006effect,kaufmann2010fluctuating,vidovic2016improving}. The gain from training on two sessions is confounded with doubling the training data. On DB5 the whole procedure ran without modification, but \rev{with 10 subjects the estimates were imprecise, and the $k = 12$ decision depended on the bootstrap seed and did not survive Holm correction, with neighbouring $k$ failing. The result neither replicates nor contradicts the GRABMyo reduction.} \rev{On putEMG, a larger cohort with a three-band, 24-electrode montage added after the main analysis, the same procedure also found a non-inferior reduction ($k^* = \PutKStarPtTwo$ of 24). Because this analysis was post hoc and the montage differs, it supports the transferability of the procedure rather than any particular retained fraction.}

\paragraph{Limitations}
The primary unit is a monopolar electrode. A conventional device measures bipolar pairs, and removing four electrodes does not remove four sensors of such a device. The exploratory bipolar re-analysis points the other way: with pairs as the unit, not even one of eight could be removed within the primary margin, although the shortfall at seven pairs was small. The 25\% reduction should therefore be read as a statement about monopolar electrodes on this array, not as evidence that a bipolar armband can shed a quarter of its sensors. The bipolar reference is itself a derived montage, and physical bipolar electrodes may behave differently. All participants were able-bodied, and all analysis was offline\rev{: filtering was zero-phase (non-causal) and windows were 250~ms, so the results are an offline benchmark and do not establish performance in real-time, causal control. Fatigue, limb posture and electrode shift were not isolated. Rankings were computed from full-array models and could shift if models were trained only on a reduced array. No imaging-based registration of electrodes to muscles was available.} The SVM was trained on a subsample and is \rev{likely} a lower bound. The H2 and H4 comparisons used a lighter comparison model than the headline curve. The count-feature threshold was a single scalar set \rev{once rather than per fold (a leave-one-out check found no effect), training-set sizes were calibrated on two held-out folds, and the bipolar threshold's derivation is not fully recorded}. $\lambda$ was not selected, and the detailed H2--H4 test families were specified after the underlying results were computed (Section~\ref{sec:deviations}). \rev{The fixed-layout analysis and the putEMG replication were post hoc, and DB5 had only 10 subjects.} Cross-model attribution, window-length sensitivity and the 17-class task were listed as exploratory in the plan and were not run.
